\section{Diagrama de caso de uso}
% Organização sugerida para uma introdução de plano de estágio

O diagrama de caso de uso é uma das partes fundamentais do projeto, pois é com ele que serão definidos os requisitos funcionais de forma clara do sistema. Na Figura~\ref{fig:caso-de-uso} são apresentados os principais atores e funcionalidades da aplicação.

\begin{figure}[H]
    \caption{Caso de uso}
    \label{fig:caso-de-uso}
    \centering
    \includegraphics[width=1\textwidth]{imagens/caso_de_uso.png}

    {\footnotesize \textbf{Fonte:} O autor.}
\end{figure}

\section{Casos de uso}

\begin{description}
    \item[\textbf{suc\_enviar\_notificacao}] Permite ao sistema enviar notificações aos usuários sobre atualizações, lembretes e eventos importantes.

    \item[\textbf{suc\_gerenciar\_autenticacao}] Permite realizar login, logout e controle de acesso dos usuários no sistema.

    \item[\textbf{suc\_gerenciar\_salas}] Permite cadastrar, editar, remover e consultar as salas disponíveis.

    \item[\textbf{suc\_gerenciar\_agendamentos}] Permite criar, atualizar, cancelar e visualizar agendamentos de salas.

    \item[\textbf{suc\_gerar\_relatorios}] Permite emitir relatórios com informações consolidadas sobre uso, agendamentos e histórico do sistema.

    \item[\textbf{suc\_gerenciar\_histórico}] Permite consultar e acompanhar o histórico de ações, reservas e alterações realizadas.

    \item[\textbf{suc\_gerenciar\_usuarios}] Permite cadastrar, editar, remover e consultar usuários do sistema.

    \item[\textbf{suc\_exportar\_PDF}] Permite exportar documentos, relatórios ou comprovantes em formato PDF.

    \item[\textbf{suc\_montar\_grade}] Permite organizar e estruturar a grade de horários ou ocupação das salas.
\end{description}
